% ============================================================================
% Fichier  : partie2/P02_C07_regles_detection.tex
% Titre    : Écrire et Éprouver une Règle de Détection
% Auteur   : MINKA MI NGUIDJOÏ Thierry Emmanuel
% ============================================================================
% Sources : MITRE ATT&CK ; pratiques d'ingénierie de détection.
% ============================================================================

\chapter{Écrire et Éprouver une Règle de Détection}
\label{chap:regles}
\chaptermark{Écrire une règle de détection}

\epigraph{%
    « Une bonne règle se juge à deux nombres~:\\
    ce qu'elle attrape, et ce qu'elle dérange pour rien. »%
}{Principe du cours}

\niveauChapitre{Intermédiaire à expert}{Chapitres~\ref{chap:siem}
et~\ref{chap:idsips}.}

\begin{boxobjectifs}
\begin{itemize}
    \item Formuler une règle de détection en français avant de la coder.
    \item Relier chaque règle à une technique ATT\&CK et à sa trace.
    \item Éprouver une règle~: vrais positifs, faux positifs, faux négatifs.
    \item Maintenir un catalogue de règles dans le temps.
    \item Rédiger la fiche d'une règle.
\end{itemize}
\end{boxobjectifs}

\bigskip

% ============================================================================
\section*{Question directrice : comment transformer « un compte de service anormal » en une alerte fiable~?}
% ============================================================================

L'idée d'une règle est simple~; la faire fonctionner sans noyer l'analyste
l'est moins. Ce chapitre apprend à écrire une règle, à l'éprouver et à la
maintenir, à partir des trois cas d'usage de l'opération BÉNOUÉ.

\medskip
\textbf{Plan du chapitre.}
\begin{enumerate}
    \item Énoncer avant de coder (section~\ref{sec:c07-enoncer}).
    \item Relier à une technique et à sa trace (section~\ref{sec:c07-attack}).
    \item Éprouver une règle (section~\ref{sec:c07-eprouver}).
    \item Maintenir le catalogue (section~\ref{sec:c07-maintenir}).
\end{enumerate}

\separateur

% ============================================================================
\section{Énoncer avant de Coder}
\label{sec:c07-enoncer}
% ============================================================================

\begin{boxprocedure}
\textbf{Énoncer une règle}
\begin{enumerate}
    \item \textbf{Quoi}~: le comportement à détecter, en une phrase.
    \item \textbf{Source}~: le ou les journaux concernés.
    \item \textbf{Condition}~: ce qui déclenche (seuil, enchaînement, contexte).
    \item \textbf{Exceptions}~: les cas légitimes à exclure.
    \item \textbf{Réponse}~: priorité, action attendue de l'analyste.
\end{enumerate}
\end{boxprocedure}

Exemple~: « Alerter si un compte de service ouvre une session interactive sur
un poste de travail (source~: journaux du domaine, événement de connexion de
type~2), hors des comptes et plages de maintenance déclarés. »

% ============================================================================
\section{Relier à une Technique et à sa Trace}
\label{sec:c07-attack}
% ============================================================================

\begin{tableaucours}{L{4.4cm}L{4.2cm}L{3.6cm}}{%
    \tableauHeader{Cas d'usage (BÉNOUÉ)} &
    \tableauHeader{Technique ATT\&CK} &
    \tableauHeader{Trace}
}
Compte de service sur un poste & Identifiants valides (usage) & Connexion du domaine \\
Création de compte hors maintenance & Persistance & Événement de création de compte \\
Service installé à distance & Mouvement latéral & Événement d'installation de service \\
\end{tableaucours}
\vspace{-0.8em}
\begin{center}{\small\itshape Tableau~7.1. Trois règles reliées à ATT\&CK~\cite{mitreattack}}\end{center}

% ============================================================================
\section{Éprouver une Règle}
\label{sec:c07-eprouver}
% ============================================================================

\begin{boxdef}
\termecle{Vrai positif}~: la règle alerte sur un vrai comportement visé.
\termecle{Faux positif}~: elle alerte à tort. \termecle{Faux négatif}~: elle
laisse passer le comportement visé. On éprouve une règle en la rejouant sur
des journaux réels et sur un jeu de test (la technique jouée en laboratoire).
\end{boxdef}

\begin{boxCRO}
\textbf{Analyse CRO d'une règle trop large.} Une règle qui collecte et
inspecte un trafic très étendu « pour ne rien rater » dégrade la
\cro{-}{+}{0}~: elle augmente la détection (fiabilité) mais expose davantage
de données (confidentialité). On cible la règle sur ce qui est nécessaire.
\end{boxCRO}

% ============================================================================
\section{Maintenir le Catalogue}
\label{sec:c07-maintenir}
% ============================================================================

\begin{boxscenario}{Opération BÉNOUÉ. Du cas d'usage à la règle éprouvée}
\textbf{Constatation.} La règle « compte de service sur un poste » alerte
24~fois en une semaine de test~: 3~vrais positifs (comptes réellement mal
utilisés), 21~faux positifs (un script de sauvegarde légitime).

\medskip
\textbf{Analyse} (confiance élevée). En excluant le compte du script de
sauvegarde déclaré, les faux positifs tombent à zéro sans perdre les vrais
positifs. Une règle n'est utile qu'une fois éprouvée et affinée.
\end{boxscenario}

\begin{boxmetier}{Analyste SOC / ingénierie de détection}
\textbf{Une règle vit.} L'ingénieur de détection versionne ses règles, note
leur technique ATT\&CK, mesure leur bruit et les révise après chaque
incident. Un catalogue de règles est un actif à entretenir, pas une
installation unique.
\end{boxmetier}

\separateur

\begin{boxresume}
\textbf{Ce qu'il faut retenir de ce chapitre~:}
\begin{itemize}
    \item \textbf{Énoncer} en français~: quoi, source, condition, exceptions,
          réponse.
    \item \textbf{Relier} chaque règle à une technique ATT\&CK et à sa trace.
    \item \textbf{Éprouver}~: vrais positifs, faux positifs, faux négatifs.
    \item \textbf{Affiner} par les exceptions, sans perdre les vrais positifs.
    \item \textbf{Maintenir} le catalogue dans le temps.
\end{itemize}
\end{boxresume}

\begin{boxinfo}
\textbf{Points de contrôle de l'élève.} Je sais~: énoncer une règle~; la
relier à ATT\&CK~; distinguer les trois cas (VP, FP, FN)~; affiner par
exceptions~; entretenir un catalogue.
\end{boxinfo}

\bigskip

\begin{boxexo}[title={\ding{45}\quad Exercice~7.1. Énoncer \niveaubase}]
Énoncez en français la règle « service installé à distance hors maintenance ».
\end{boxexo}

\begin{boxexo}[title={\ding{45}\quad Exercice~7.2. Trois cas \niveaubase}]
Définissez vrai positif, faux positif, faux négatif.
\end{boxexo}

\begin{boxexo}[title={\ding{45}\quad Exercice~7.3. Technique \niveaubase}]
À quelle tactique ATT\&CK relier « création de compte hors maintenance »~?
\end{boxexo}

\begin{boxexo}[title={\ding{45}\quad Exercice~7.4. Affiner \niveaumoyen}]
Comment ramener à zéro les 21~faux positifs sans perdre les 3~vrais positifs~?
\end{boxexo}

\begin{boxexo}[title={\ding{45}\quad Exercice~7.5. Faux négatif \niveaumoyen}]
Pourquoi un faux négatif est-il plus grave qu'un faux positif~?
\end{boxexo}

\begin{boxexo}[title={\ding{45}\quad Exercice~7.6. CRO \niveaumoyen}]
Expliquez l'analyse CRO d'une règle trop large.
\end{boxexo}

\begin{boxexo}[title={\ding{45}\quad Exercice~7.7. Jeu de test \niveaumoyen}]
Comment constituer un jeu de test pour éprouver une règle~?
\end{boxexo}

\begin{boxexo}[title={\ding{45}\quad Exercice~7.8. Fiche de règle \niveaumoyen}]
Rédigez la fiche complète d'une des trois règles de BÉNOUÉ.
\end{boxexo}

\begin{boxexo}[title={\ding{45}\quad Exercice~7.9. Catalogue \niveauexpert}]
Organisez un catalogue de règles~: versionnage, métadonnées, révision.
\end{boxexo}

\begin{boxexo}[title={\ding{45}\quad Exercice~7.10. Couverture \niveauexpert}]
À partir d'ATT\&CK, proposez cinq règles qui augmenteraient la couverture de
détection de LiZbiZ.
\end{boxexo}

\bigskip

\begin{boxinfo}
\textbf{Pour aller plus loin.}
\begin{itemize}
    \item MITRE ATT\&CK~\cite{mitreattack}~; NIST SP~800-61r2~\cite{nist80061r2}.
\end{itemize}
\end{boxinfo}

\begin{boxinfo}
\textbf{Transition.} La Partie~III pose le cadre normatif et juridique de
tout ce qui précède~: le Chapitre~\ref{chap:normes} ouvre sur les normes de
l'incident et de la preuve.
\end{boxinfo}
